\section*{Abstract}

Unit-test generation with large language models (LLMs) is becoming increasingly common,
but most existing studies still rely on simple, curated benchmarks that neither reflect
the complexity of real-world code nor control for it. This study tests the one-shot
capability of \texttt{gpt-4o-mini} on 120 real-world functions (60 Java, 60 Python) of
medium cyclomatic complexity (5--10), mined from ten open-source repositories at pinned
commits and evaluated directly inside their original projects. Its results are compared
against EvoSuite and Randoop (Java) and Pynguin (Python) on branch coverage and mutation
score, using non-parametric tests ($\alpha = 0.05$). Across all 120 functions, median
coverage and median mutation score are both 0\%. Only 28/120 suites (23.3\%) compile and
actually test the intended function; their median coverage is 75.0\%, still below the
80\% bar. Over suites with a measurable mutation score, the median is 36.84\% for Python
and 0\% for Java, both under the 60\% threshold. Against the baselines, \texttt{gpt-4o-mini} is beaten outright by both
Java tools (EvoSuite and Randoop, $r = -1.00$ for each). It beats Pynguin only because
Pynguin failed completely on this dataset (median coverage 0.0\%). Within the CC~5--10
band there is no clear link between complexity and test quality; what decides the
outcome is whether a suite compiles and tests the right function. Python failures are
mostly bad imports, while Java calls the wrong method without any error --- something
only visible when running on real project code. Overall, one-shot LLM generation,
without retrieval or repair, cannot yet replace automated test-generation tools on real
code of medium complexity. The results support adding feedback loops and richer context
when generating tests.
